% 	Name		:: 	sthlm Beamer Theme  HEAVILY based on the hsrmbeamer theme (Benjamin Weiss)
%	Author		:: 	Mark Hendry Olson (mark@hendryolson.com)
%	Created		::	2013-07-31
%	Updated		::	June 18, 2015 at 08:45
%	Version		:: 	1.0.2
%	Email		:: 	hendryolson@gmail.com
%	Website		:: 	http://v42.com
%
% 	License		:: 	This file may be distributed and/or modified under the
%                  	GNU Public License.
%
%	Description	::	This presentation is a demonstration of the sthlm beamer
%					theme, which is HEAVILY based on the HSRM beamer theme created by Benjamin Weiss
%					(benjamin.weiss@student.hs-rm.de), which can be found on GitHub
%					<https://github.com/hsrmbeamertheme/hsrmbeamertheme>.


%-=-=-=-=-=-=-=-=-=-=-=-=-=-=-=-=-=-=-=-=-=-=-=-=
%
%        LOADING DOCUMENT
%
%-=-=-=-=-=-=-=-=-=-=-=-=-=-=-=-=-=-=-=-=-=-=-=-=

\documentclass[newPxFont]{beamer}
\usetheme{sthlm}
%\usecolortheme{sthlmv42}

%-=-=-=-=-=-=-=-=-=-=-=-=-=-=-=-=-=-=-=-=-=-=-=-=
%        LOADING PACKAGES
%-=-=-=-=-=-=-=-=-=-=-=-=-=-=-=-=-=-=-=-=-=-=-=-=
\usepackage[utf8]{inputenc}
\usepackage[T1]{fontenc}

%\usepackage{chronology}
\usepackage{chronosys}
\usepackage{subfigure}

\newcommand{\tabitem}{%
  \usebeamertemplate{itemize item}\hspace*{\labelsep}}

%\renewcommand{\event}[3][e]{%
%  \pgfmathsetlength\xstop{(#2-\theyearstart)*\unit}%
%  \ifx #1e%
%    \draw[fill=black,draw=none,opacity=0.5]%
%      (\xstop, 0) circle (.2\unit)%
%      node[opacity=1,rotate=45,right=.2\unit] {#3};%
%  \else%
%    \pgfmathsetlength\xstart{(#1-\theyearstart)*\unit}%
%    \draw[fill=black,draw=none,opacity=0.5,rounded corners=.1\unit]%
%      (\xstart,-.1\unit) rectangle%
%      node[opacity=1,rotate=45,right=.2\unit] {#3} (\xstop,.1\unit);%
%  \fi}%

%-=-=-=-=-=-=-=-=-=-=-=-=-=-=-=-=-=-=-=-=-=-=-=-=
%        BEAMER OPTIONS
%-=-=-=-=-=-=-=-=-=-=-=-=-=-=-=-=-=-=-=-=-=-=-=-=

%\setbeameroption{show notes}

%-=-=-=-=-=-=-=-=-=-=-=-=-=-=-=-=-=-=-=-=-=-=-=-=
%
%	PRESENTATION INFORMATION
%
%-=-=-=-=-=-=-=-=-=-=-=-=-=-=-=-=-=-=-=-=-=-=-=-=

\title{The tree perceptions in a Sahelian context}
\subtitle{ANR « FUTURE SAHEL »}
%\date{\small{\jobname}}
%\date{\today}
\date{September 2017}
\author{\texttt{E. Delay,  K. Niang, A. Ka, J.-L. Peiry, D. Goffner}}
\institute{Presentation at Stockholm Resilience Center}

\hypersetup{
pdfauthor = {E. DELAY},
pdfsubject = {Presentation at SRC},
pdfkeywords = {Great Green Wall, Senegal, OHM},
pdfmoddate= {D:\pdfdate},
pdfcreator = {pdflatex}
}

\begin{document}

%-=-=-=-=-=-=-=-=-=-=-=-=-=-=-=-=-=-=-=-=-=-=-=-=
%
%	TITLE PAGE
%
%-=-=-=-=-=-=-=-=-=-=-=-=-=-=-=-=-=-=-=-=-=-=-=-=

\maketitle

%\begin{frame}[plain]
%	\titlepage
%\end{frame}

%-=-=-=-=-=-=-=-=-=-=-=-=-=-=-=-=-=-=-=-=-=-=-=-=
%
%	TABLE OF CONTENTS: OVERVIEW
%
%-=-=-=-=-=-=-=-=-=-=-=-=-=-=-=-=-=-=-=-=-=-=-=-=
%\section*{Overview}
%\begin{frame}{Overview}
%% For longer presentations use hideallsubsections option
%\tableofcontents[hideallsubsections]
%\end{frame}

%-=-=-=-=-=-=-=-=-=-=-=-=-=-=-=-=-=-=-=-=-=-=-=-=
%	FRAME: INTRODUCTION
%-=-=-=-=-=-=-=-=-=-=-=-=-=-=-=-=-=-=-=-=-=-=-=-=

\section{Introduction}

%-=-=-=-=-=-=-=-=-=-=-=-=-=-=-=-=-=-=-=-=-=-=-=-=
%	FRAME: CONTEXT
%-=-=-=-=-=-=-=-=-=-=-=-=-=-=-=-=-=-=-=-=-=-=-=-=
\begin{frame}[c]{Context}
\vspace{-1cm}
% Questionnement autour de la place de l'arbre
\begin{itemize}
  \item \textbf{Observation} : populations mainly focused on pastoralism or agro-pastoralism.
  \item \textbf{Objective} : understand the "traditional" place of ESS products from trees to increase socio-environmental resilience.
\end{itemize}

What are the uses of trees as providers of resources and vectors of ecosystem services (medicinal, cosmetics, food, fodder, etc.)?

\begin{figure}
	\centering
	\includegraphics[width = 0.5\textwidth]{img/2007_Podor.JPG}
	\caption{FutureSahel Work Packages}
\end{figure}
\end{frame}

%-=-=-=-=-=-=-=-=-=-=-=-=-=-=-=-=-=-=-=-=-=-=-=-=
%	FRAME: Résilience et stockolm
%-=-=-=-=-=-=-=-=-=-=-=-=-=-=-=-=-=-=-=-=-=-=-=-=
\begin{frame}[c]{Socio-environmental resilience: a practical approach}
\vspace{-1cm}
Can we consider the trees eco-systemic services in the Sahelian context as a factor of resilience for the populations and the environment?


\begin{figure}
	\centering
	\includegraphics[width = 0.65\textwidth]{img/PA310152.JPG}
	\caption{a tree eco-systemic services}
\end{figure}
\end{frame}


\section{methodology - ethnobotanical entry}

%-=-=-=-=-=-=-=-=-=-=-=-=-=-=-=-=-=-=-=-=-=-=-=-=
%	FRAME: ethnobotanique
%-=-=-=-=-=-=-=-=-=-=-=-=-=-=-=-=-=-=-=-=-=-=-=-=
\begin{frame}[c]{Ethnobotanical interviews}
\vspace{-1cm}

Structured semi-structured interviews with stakeholders (individual and focus groups) in order to understand the problems and realities of the field, to identify the essential elements to be taken into account.

Know/understand the uses, the rising strategies, the seasonality, the methods of valorisations.

%acteurs photos jean-luc

\begin{figure}
	\hspace{-1.5em}
	\includegraphics[height = 3.5cm]{img/group3.png}~
  \includegraphics[height = 3.5cm]{img/group2.png}
\end{figure}
\end{frame}

%-=-=-=-=-=-=-=-=-=-=-=-=-=-=-=-=-=-=-=-=-=-=-=-=
%	FRAME: Questionnaires
%-=-=-=-=-=-=-=-=-=-=-=-=-=-=-=-=-=-=-=-=-=-=-=-=
\begin{frame}[c]{The questionnaire approach}
\vspace{-1cm}
A specific questionnaire is used to quantify the information already obtained during the interviews.

Completion of 30-40 questionnaires per site (a total of 200 questionnaires for all sites).
%Image googleEarth des fenetres
\begin{figure}
	\centering
	\includegraphics[width = 0.6\textwidth]{img/Khoudia.png}
\end{figure}

\end{frame}

%-=-=-=-=-=-=-=-=-=-=-=-=-=-=-=-=-=-=-=-=-=-=-=-=
%	FRAME: Ecologie du paysage
%-=-=-=-=-=-=-=-=-=-=-=-=-=-=-=-=-=-=-=-=-=-=-=-=
\section{methodology - contributions of landscape ecology}

\begin{frame}[c]{The landscape ecology approach 1/3}
\vspace{-1cm}
A work inspired by the work proposed by Sinare \textit{et al.} (2016)
\begin{figure}
	\centering
	\includegraphics[height = 3.5cm]{img/ggearth}
  \includegraphics[height = 3.5cm]{img/Sinare_et_al2016}
  \caption{\small{Screenshot Google-earth and maps from Sinare \textit{et al.} (2016)}}
\end{figure}
\end{frame}

%-=-=-=-=-=-=-=-=-=-=-=-=-=-=-=-=-=-=-=-=-=-=-=-=
%	FRAME: Unité paysagères
%-=-=-=-=-=-=-=-=-=-=-=-=-=-=-=-=-=-=-=-=-=-=-=-=

\begin{frame}[c]{Landscape ecology in workshop 2/3}
\vspace{-1cm}
\begin{figure}
  \subfigure[]{\includegraphics[height = 2.5cm]{img/photo_mission/01_Plage_filao}}
  \subfigure[]{\includegraphics[height = 2.5cm]{img/photo_mission/02_Champs_irrigue}}\\
  \subfigure[]{\includegraphics[height = 2.5cm]{img/photo_mission/03_Brousse_peu_dense}}
  \subfigure[]{\includegraphics[height = 2.5cm]{img/photo_mission/05_Brousse_dense}}
\end{figure}

\end{frame}

\begin{frame}[c]{Landscape ecology in workshop 3/3}
\vspace{-1cm}
\begin{figure}
	\subfigure[]{\includegraphics[height = 2.5cm]{img/photo_mission/06_Mares}}
  \subfigure[]{\includegraphics[height = 2.5cm]{img/photo_mission/07_Homestead}}\\
  \subfigure[]{\includegraphics[height = 2.5cm]{img/photo_mission/08_Agriculture_pluvial}}
  \subfigure[]{\includegraphics[height = 2.5cm]{img/photo_mission/10_Parcelle_reboisement}}
\end{figure}
\end{frame}


%-=-=-=-=-=-=-=-=-=-=-=-=-=-=-=-=-=-=-=-=-=-=-=-=
%	FRAME: Participation
%-=-=-=-=-=-=-=-=-=-=-=-=-=-=-=-=-=-=-=-=-=-=-=-=
\section{methodology - the participatory approach}
\begin{frame}[c]{Workshop}
\vspace{-1cm}
Workshops with stockholders with objective to highlight actors importance of trees
in each structure of the landscape.

\begin{itemize}
  \item bring space structure out (orientation + mark in the living space),
  \item identify the spatial units in the region,
  \item quantify the use of the trees in these units.
\end{itemize}
\begin{figure}
	\centering
	\includegraphics[width = 0.6\textwidth,angle=180]{img/DSC_1789}
\end{figure}
\end{frame}


%-=-=-=-=-=-=-=-=-=-=-=-=-=-=-=-=-=-=-=-=-=-=-=-=
%	FRAME: Quantification
%-=-=-=-=-=-=-=-=-=-=-=-=-=-=-=-=-=-=-=-=-=-=-=-=

\begin{frame}[c]{Workshop quantification}
\vspace{-1cm}
Quantify the use of the trees in these units using mind maps and glass beads.

\begin{figure}
  \centering
  \includegraphics[width = 0.8\textwidth]{img/DSC_1793}
\end{figure}

\end{frame}

%-=-=-=-=-=-=-=-=-=-=-=-=-=-=-=-=-=-=-=-=-=-=-=-=
%	FRAME: Résultats
%-=-=-=-=-=-=-=-=-=-=-=-=-=-=-=-=-=-=-=-=-=-=-=-=

\section{firsts results}

\begin{frame}[c]{Result 1 : identification of the living space}
\vspace{-1cm}
\begin{figure}
	\subfigure[]{\includegraphics[height = 4.5cm]{img/DSC_0371}}
  \subfigure[]{\includegraphics[height = 4.5cm]{img/map_louguere.png}}
\end{figure}
\end{frame}

%-=-=-=-=-=-=-=-=-=-=-=-=-=-=-=-=-=-=-=-=-=-=-=-=
%	FRAME: Classification
%-=-=-=-=-=-=-=-=-=-=-=-=-=-=-=-=-=-=-=-=-=-=-=-=
\begin{frame}[c]{Result 2 : patches classification}
\vspace{-1cm}
\begin{figure}
  \centering
  \includegraphics[width = \textwidth]{img/process}
\end{figure}
\end{frame}

\begin{frame}[c]{Result 2 : patches classification}
\vspace{-1cm}
\begin{figure}
  \hspace*{-1cm}
	\subfigure[]{\includegraphics[height = 4.7cm]{img/Classif}}~
  \subfigure[]{\includegraphics[height = 4.7cm]{img/camp_agri}}~
  \subfigure[]{\includegraphics[height = 4.7cm]{img/distance}}
\end{figure}
\end{frame}

%-=-=-=-=-=-=-=-=-=-=-=-=-=-=-=-=-=-=-=-=-=-=-=-=
%	FRAME: Quantification
%-=-=-=-=-=-=-=-=-=-=-=-=-=-=-=-=-=-=-=-=-=-=-=-=
\begin{frame}[c]{Result 3 : Quantification}
\vspace{-1cm}
\begin{figure}
  \centering
  \includegraphics[width = 0.8\textwidth]{img/papier-figBar1}
  \caption{Trees use by spatial units}
\end{figure}

\end{frame}

\begin{frame}[c]{Result 3 : Quantification}
\vspace{-1cm}
\begin{figure}
  \centering
  \includegraphics[width = \textwidth]{img/papier-figBar2}
  \caption{Usage de l'arbre par unité spatiale}
\end{figure}

\end{frame}

\begin{frame}[c]{Result 4 : spatiale quantification}
\vspace{-1cm}
\begin{figure}
  \centering
  \includegraphics[width = 0.8\textwidth]{img/usages_sp_distances}
  \caption{sp}
\end{figure}

\end{frame}\begin{frame}[c]{Result 4 : spatiale quantification}
\vspace{-1cm}
\begin{figure}
  \centering
  \includegraphics[width = 0.7\textwidth]{img/sum_use_sp}
  \caption{Sum of trees uses}
\end{figure}

\end{frame}

%-=-=-=-=-=-=-=-=-=-=-=-=-=-=-=-=-=-=-=-=-=-=-=-=
%	FRAME: Conclusion
%-=-=-=-=-=-=-=-=-=-=-=-=-=-=-=-=-=-=-=-=-=-=-=-=
\section{Conclusion}
%-=-=-=-=-=-=-=-=-=-=-=-=-=-=-=-=-=-=-=-=-=-=-=-=
%	FRAME: Intégration à ComMod
%-=-=-=-=-=-=-=-=-=-=-=-=-=-=-=-=-=-=-=-=-=-=-=-=
\begin{frame}[c]{From trees to ComMod}
\vspace{-1cm}
An interdisciplinary approach:
\begin{itemize}
  \item ethnobotanical  investigation work makes it possible to identify socio-spatial practices,
  \item landscape ecology, allow spatial identification of places contribution for those practices
  \item The formalisation of practices and the identification of place leads to the formalisation of behavioural models and therefore to a prospective approach validated by the actors.
\end{itemize}
\end{frame}


%-=-=-=-=-=-=-=-=-=-=-=-=-=-=-=-=-=-=-=-=-=-=-=-=
%	FRAME: MERCI DE VOTRE ATTENTION
%-=-=-=-=-=-=-=-=-=-=-=-=-=-=-=-=-=-=-=-=-=-=-=-=
{
\usebackgroundtemplate{
\includegraphics[width=\paperwidth]{img/podor_fin}}%
\begin{frame}
  \vspace{-1em}
  \begin{minipage}[t][.8\textheight]{\textwidth}
    \color{\cnGrey}{\LARGE{Thank you for your attention}}

    \vfill

  %\hfill \small{Photo credit : Thomas m-louis. sur \includegraphics[height=0.55cm]{img/flickr_logo}}
  \end{minipage}
  \vspace{-3.5em}
  \centering
	You can find this presentation on github\includegraphics[height=0.85cm]{img/github}

\end{frame}
}


\end{document}
